\begin{flushleft}

	\subsection*{Abstract}
		%\vspace{1.5cm}

		\subsubsection*{The AirScout project:}

		The main goal of the project AirScout is the development of a web application
		for the visualization and analysis of environmental data captured with mobile
		sensor boxes. These sensor boxes will be developed by the electrical
		engineering department and mounted on bicycles to collect data across St. Pölten and the surrounding areas. The web application will graphically and
		interactively present this sensor data on a map and a leaderboard, allowing
		users to explore and analyze the collected information in an engaging way.

		A legacy of the project already exists, but had major performance issues,
		leaked features, and was not designed for scalability. The goal of this project
		is to create a new, modern, and efficient implementation of the AirScout
		platform that addresses these issues and provides a solid foundation for future
		development and expansion.

		\subsubsection*{Overview on the diploma thesis:}

		This diploma thesis deals with the technical foundations and tech-stack
		decisions for the AirScout project. The structure is divided into the following
		key areas:

		\begin{itemize}
			\item \textbf{Data Collection and Preparation} (Luca Reisenbichler, Chapter~\ref{chap:luca-introduction}): Examination of sensor functionality, measurement collection, and raw data preparation. Covers data sources (Chapter~\ref{title:data-sources}) and data processing (Chapter~\ref{title:processing}) for handling sensor noise, missing values, and outliers.

			\item \textbf{Time-Series Databases} (Benjamin Friedl, Chapter~\ref{chap:benjamin-introduction}): Focus on efficient storage and querying of long-term data, including performance optimization, benchmarking, and database system evaluation for the AirScout use case.

			\item \textbf{Spatial Visualization} (Christian Patzl, Chapter~\ref{chap:christian-introduction}): Theoretical foundations and implementation of spatial visualization in AirScout, enabling users to inspect sensor values across space and time.

			\item \textbf{3D Representation} (Tobias Eisinger, Chapter~\ref{chap:tobias-introduction}): Mathematical and technical foundations for 3D visualization of sensor locations, providing intuitive geographic context and better understanding of physical sensor placement.

			\item \textbf{Mobile Access} (Valentin Hacker, Chapter~\ref{chap:why_needed}): Evaluation of cross-platform mobile frameworks and implementation of a mobile application for convenient access to measurement data and device management.
		\end{itemize}
\end{flushleft}